% Parent-ready source/status fragment. No preamble.
\section{Source audit, integration decisions, and further research}
\label{status:section}

The source anchor for this continuation is ProveIt revision
\texttt{03112a7d3b91} \cite{ProveIt}.
The audit concerns the incoming passages used in the proofs and the
manuscript's relevant status records. It is not a certification of
every result in the repository. The three main advances have the
following precise relation to that source state.

\subsection{What the present results change}

\begin{enumerate}
\item \textbf{Joint Gamma--polygamma moments.}
The Gamma-normalization report's
\path{sections/02_gaussian.tex} proves the paired deformation and
the single-trigamma Gaussian moments through degree $n+1$;
research question Q2 in
\path{sections/06_audit_research.tex} asks for the joint
polygamma moments \cite{GammaIncoming}.
Theorem~\ref{gc:main} separates all joint cumulants through total
weight $2n+2$, with weight $q+1$ assigned to $R_q$.
It includes digamma and even polygamma orders. Formula
\eqref{gc:moment} then evaluates every retained ordinary moment,
and Corollary~\ref{gc:wick} gives the centered positive-order Wick
rule. This answers the joint-separation question in the stated
weighted range; the beyond-range classification remains open.
The nonzero small-$\sigma$ skewness limit in
Proposition~\ref{gc:sharpness} supplies an obstruction relevant to
the same report's Q1, but does not evaluate the general first
excluded trigamma moment.
The underlying Gustafson integral remains classical
\cite{DerkachovManashov}; its Gamma normalizer and the paired law
retain their incoming attribution.

\item \textbf{Effective reconstruction and the complete gap-eight ray.}
Section~8 of the ordered-boundary report gives eventual
nonsingularity at each fixed gap; Section~9, research question~1,
asks for an effective threshold and finite completion
\cite{OrderedIncoming}. Theorem~\ref{eg:effective} supplies an
explicit computable $B(m)$ for every fixed half-gap $m$.
The later gap-six report's
\path{sections/05_audit_and_questions.tex}, question Q1,
explicitly asks for complete gap-eight coverage \cite{GapSix}.
Theorem~\ref{eg:eight-main} proves this for every
$(p,q)=(2b-8,2b)$, $b\ge5$, using a proved analytic tail and an
exact finite-field completion. Corollary~\ref{eg:eight-identities}
turns that result into ordinary convergent polygamma-product
identities. The unrestricted all-pairs prefix conjecture remains
open. The finite certificate proves nonvanishing, and does not
prove positivity at every intervening finite index.

\item \textbf{Stieltjes collisions at different monomial scales.}
The gap-six report's Section~4 treats fixed positive rates
$a=\lambda\epsilon$, $b=\mu\epsilon$; its question Q10 asks for
$a=\lambda\epsilon^\alpha$, $b=\mu\epsilon^\beta$ with
$\alpha\ne\beta$ \cite{GapSix}.
Theorem~\ref{mc:main} evaluates the constant for every positive
monomial pair and every fixed pair of Stieltjes indices and
argument derivatives. Its proof also specifies all divergent
terms through that constant by finite truncation. The three
resonance conditions in \eqref{mc:three-mechanisms} identify the
additional contacts, including the singular--singular contact.
The earlier lattice-contact report's Q9 asked only for fixed
unequal rates \cite{LatticeIncoming}; that problem was already
answered by the gap-six report and is not reclaimed here.
General nonmonomial paths and three independently shifted factors
remain outside the present theorem. The trigamma-tail contact
family \eqref{mc:stieltjes-tail-contact} and its polylogarithmic
generator and primitive are explicit evaluated consequences of
the new scale classification.
\end{enumerate}

The Barnes-integral companion supplies additional ordinary integral
evaluations obtained by explicitly differentiating a classical Barnes
lemma \cite{DLMF}. It should be credited in that form. A bounded
literature search and comparison with the inspected incoming material
support the stated continuation choices; they do not establish global
historical novelty.

\subsection{Errors, status corrections, and essential qualifications}

No newly false inherited mathematical theorem was found in the
specific dependencies used for these results. The changes above are
proved extensions and research-status updates. They should not be
presented as corrections to formulas which the sources stated correctly.
In particular, the incoming Gamma report already restricts its finite
Gaussian law, and the gap-six report already distinguishes finite-field
nonvanishing from a real-sign theorem.

One source-list reconciliation is necessary. The gap-six report's Q4
and the ordered-boundary report's question~2 ask for evaluations of
$K_m$. The later Gamma-normalization report already supplies them
\cite{GammaIncoming,GapSix,OrderedIncoming}. These questions should be
marked answered by that predecessor, rather than counted among the
present results. The repeated README paragraph about the 1,032-input
first-letter-$X$ family, rank 1,008 and kernel dimension 24 remains
present at the pinned revision. Earlier incoming reports already
identified this editorial duplication. Removing its second occurrence
would be an inherited editorial cleanup, not a new mathematical erratum.

Three conventions are material to the proofs. The Gamma moment
classification retains its total-weight restriction: beyond that range,
Newton relations prevent unrestricted comparison of power-sum
coefficients. The Bernoulli reconstruction matrix retains the omission
of a lower binomial index zero; evaluating that entry as
$\binom n0=1$ changes the problem. The Stieltjes constant is taken in
the displayed small parameter, with the fixed-coordinate finite part
kept distinct from the path-dependent constant. Spectral pole pairs
must be combined before their removable values and derivatives are
taken. These are integration requirements, not newly discovered
mistakes in the cited sources.

The manuscript's short $S_6$ and revised $S_8$ candidates remain open;
the old rejected $S_8$ candidate remains rejected. None of the Gamma
moment, reconstruction, or collision theorems implies those polylogarithmic
period identities. Numerical proximity and rank in a specified formal
relation space do not supply their missing exact equalities.
The verification artifacts here likewise distinguish exact finite
certificates from floating-point diagnostics. The analytic arguments
justify convergence, continuation and coefficient extraction.

\subsection{Further questions directed toward exact identities}

\begin{enumerate}
\item \textbf{Evaluate the first excluded polygamma cumulants.}
At total weight $2n+3$ and above, determine the relations induced
by $e_{2n+3},e_{2n+4},\ldots$ and identify additional integral
identities that separate the remaining cumulants. The immediate
scalar target is the exact $(n+2)$nd trigamma Gaussian defect at
general $n,\sigma$. The proved $n=1$ small-shift limit is a
normalization constraint on any proposed formula.

\item \textbf{Develop joint moment formulas for unequal Gamma clusters.}
Allow several distinct base shifts with specified multiplicities.
Determine which mixed polygamma moments can be separated by
independent deformations within each cluster, how their finite
coordinate relations change when clusters merge, and which new
ordinary integrals arise at the first unresolved degrees. A useful
answer would give explicit polygamma or finite hypergeometric
coefficients, rather than only a differentiated generating integral.

\item \textbf{Find closed reconstruction rows and identities in $b$.}
The gap-eight theorem specifies $P_b$ by a rational inverse matrix.
Find finite-sum descriptions or recurrences for its coefficients and
for the corresponding subtraction polynomial. Decide which inverse
rows yield especially short zeta-valued identities. The exact
leading coefficient and the ordinary sum in
\eqref{eg:eight-sum} provide tests for every proposed recurrence.

\item \textbf{Replace fixed-gap finite completion by structural algebra.}
Find a recurrence invariant, factorization, or appropriate
variation-diminishing statement that proves nonvanishing for an
entire gap ray without checking every exceptional index separately.
Such a theorem must respect the known early sign behavior.
Connecting this structure with the fixed-smaller-order results is a
concrete route toward the unrestricted unequal-even prefix problem.

\item \textbf{Compute collision constants for nonmonomial paths.}
Take $a=\lambda\epsilon^\alpha A(\epsilon)$ and
$b=\mu\epsilon^\beta B(\epsilon)$, with prescribed analytic or
power--logarithm expansions and nonzero leading coefficients.
Determine the finite jets of $A,B$ that can affect a specified
constant and derive its exact transformation law under a change of
the small parameter. The monomial resonance rules identify the
base exponents with which these additional coefficients interact.

\item \textbf{Evaluate three-factor Stieltjes collision contacts.}
For three independently shifted derivatives, classify contributions
from every singular subset and pair all intersecting spectral poles
before differentiating. Seek a finite evaluated contact formula in
logarithms, generalized harmonic numbers and zeta derivatives, and
identify precisely where multivariable polylogarithms or other new
functions become necessary.

\item \textbf{Generate and integrate the higher contact families.}
For fixed Stieltjes and derivative orders, sum the integer-resonance
contacts against $z^j$ and determine their polylogarithmic alphabet.
The trigamma-tail generator and its anchored primitive are the first
case. Compute further Euler primitives, ordinary or finite-part
endpoint values, and character-weighted versions; determine which
reduce to zeta products and which retain genuine multiple
polylogarithms.

\item \textbf{Construct an explicit functional route to $S_6$ or revised
$S_8$.}
If one of the new generators or integral deformations is to inform
these period problems, first exhibit its exact map to the canonical
residual, including all products and branch constants. A new
functional or iterated-integral identity outside the tested relation
spaces would constitute progress; an unrelated rank theorem or
additional decimal agreement would not.
\end{enumerate}
